\section{Setup}\label{sec:sp:setup}

\subsection{Tastes and technologies}
\label{sec:sp:setup:primitives}


\smalltitle{Preferences.} A representative family dynasty with an infinite horizon derives utility from consumption and disutility from the stock of pollution. Lifetime utility at time zero is
\begin{align}
U_0 &= \int_0^{\infty}\left[u(C)-v(P)\right]e^{-\rho t}\,dt,
&&
u'>0,\quad u''<0,\quad v'>0,\quad v''>0,\quad \rho>0.
\label{eq:sp:setup:utility}
\end{align}

\smalltitle{Final goods.} Output is produced from capital and raw materials, is reduced by pollution, and is subject to exogenous technological progress, represented by the explicit $t$ argument. We include a minimum material requirement $\bar{R}\geq0$ , below which the technology is not defined:\footnote{We read $\bar{R}\geq0$ as a reduced form for the irreducible physical substrate of production, which a one-material model can only represent as a quantity floor. We define it in levels thus allowing for \emph{relative} but not \emph{absolute} dematerialization.}
\begin{align}
Y &= F\left(K^Y,R,P,t\right)\;\equiv\;
\begin{cases}
\mathcal F\left(K^Y,R^{e},P,t\right), & R\ge\bar R,
\\[0.2em]
0, & R<\bar R,
\end{cases}
&
R^{e} &\equiv R-\bar R,
\label{eq:sp:setup:output}
\end{align}
with the usual properties imposed on the effective argument,
\begin{align}
\mathcal F_K>0,\quad \mathcal F_{KK}<0,\quad \mathcal F_{R^e}>0,\quad \mathcal F_{R^eR^e}<0,\quad \mathcal F_P<0,\quad \mathcal F_t\ge0 .
\label{eq:sp:setup:output-props}
\end{align}
Setting $\bar R=0$ returns the standard specification $Y=\mathcal F\left(K^Y,R,P,t\right)$, and every statement below that does not mention $\bar R$ explicitly holds for both cases.


\smalltitle{The waste stock.} 
Physical waste flows $W$ accumulates in a stock $\mathcal{W}$ of waste held within the economy.\footnote{We think about landfills, stockpiles, discarded durables etc.. In the discrete-time version of the model, we avoid this by simply assuming that the recycling process returns recycled outputs from waste at $t+1$ instead.} Material is drawn out of the stock for handling at rate $\mu>0$:
\begin{align}
\dot{\mathcal W} &= W-\mu\mathcal W,
&
H &\equiv \mu\mathcal W,
&&
\mathcal W(0)\text{ given},
\label{eq:sp:setup:wastestock}
\end{align}
where $H$ is the \emph{handled flow}: the tonnage collected and processed per unit of time. Everything downstream of generation is levied on $H$ rather than on $W$. We treat the handling rate $\mu$ as a technological constant which prevents recycling from occurring instantaneous and avoids indeterminacy/degenerate solutions in the long run.


\smalltitle{Materials and recycling.} Raw material input is the sum of virgin and recycled materials,
\begin{align}
R &= N + R^R,
\label{eq:sp:setup:materials}
\end{align}
and recycled materials are produced from the treated share $\varpi H$ of handled waste using capital $K^R$:
\begin{align}
R^R &= a\!\left(x,t\right)\varpi H,
\qquad
x\equiv\frac{K^R}{\varpi H},
&&
a(0,t)=0,\quad a_x>0,\quad a_{xx}<0,\quad
\lim_{x\to\infty}a(x,t)=\bar a\in\left(0,1\right].
\label{eq:sp:setup:recycling}
\end{align}
Again, we include the exogenous time element to represent technological change assuming that the marginal change over time is non-negative for all $x$, i.e. $a_t\geq0$, but the boundaries remain the same.

The upper bound $\bar{a}$ is the \emph{yield ceiling} and has important implications for the long run. We distinguish between the following two scenarios:
\begin{itemize}
\item[(H)] the \emph{hard ceiling}, $\bar a<1$: however much capital is devoted to recycling, an irreducible fraction $1-\bar a$ of the treated material is lost to dissipation, contamination and quality degradation. 
\item[(S)] the \emph{soft ceiling}, $\bar a=1$: perfect recycling is priced rather than forbidden. Yield one is approached, but only as capital per treated tonne diverges, so the cost structure alone decides how circular the economy becomes.
\end{itemize}
The distinction is an assumption about matter, not an implication of thermodynamics, and it is worth being precise about which is which.\footnote{The second law prices re-concentration but does not forbid it: the minimum work needed to separate a target material from a collected stream is the Gibbs free energy of mixing, finite for any finite dilution, and the economy is an open system fed by low-entropy solar radiation. The stronger claim that complete recycling of matter is physically impossible (Georgescu-Roegen's proposed ``fourth law'' of thermodynamics \parencite{georgescu1971,georgescu1977}) is generally rejected \parencite{bianciardi1993,ayres1999}. The entropic content of the technology therefore lives in the \emph{curvature} of $a\left(\cdot\right)$, which is imposed in both cases: recovering the marginal tonne from an ever leaner residue requires ever more capital per treated tonne. The hard ceiling is strictly stronger (even unbounded input cannot close the loop) and we read it as a reduced form for two features a one-material model cannot carry explicitly: quality degradation (contamination, downcycling), which with a single homogeneous material can only be represented as quantity loss, and dissipative losses at dilutions where recovery, while physically possible, is economically out of reach at any conceivable price \parencite{reck2012,cullen2017}. Case~(S) is the entropic-cost-only benchmark in which neither reduced form is invoked.} 


Write $\varepsilon\equiv a'x/a$ for the elasticity of the yield function. It is convenient to have this technology available in constant-returns form. Let $T\equiv\varpi H$ denote treated waste and write
\begin{align}
R^R &= \mathcal R\left(T,K^R,t\right)\;\equiv\;a\!\left(\frac{K^R}{T},t\right)T,
&
x&=\frac{K^R}{T},
\label{eq:sp:setup:recycling-crs}
\end{align}
which is~\eqref{eq:sp:setup:recycling} rewritten. Because $a$ depends on $K^R$ and $T$ only through
their ratio, $\mathcal R$ is homogeneous of degree one, and its two partial derivatives are
\begin{align}
\mathcal R_T &= a\left(x\right)-a'\!\left(x\right)x \;\equiv\;\alpha,
&
\mathcal R_{K^R} &= a'\!\left(x\right),
&
\mathcal R &= \alpha T+a'K^R .
\label{eq:sp:setup:recycling-derivatives}
\end{align}
We refer to $\alpha=a\left(1-\varepsilon\right)$ as the \emph{marginal recycling yield}: the extra
tonne of secondary material obtained from one extra tonne of treated waste, holding recycling
capital fixed. It is strictly below the average yield $a$ whenever $\varepsilon>0$.\footnote{We note that $\alpha>0$ follows without further assumptions here: Concavity gives
$a\left(0\right)\le a\left(x\right)-a'\left(x\right)x$ for every $x$, so $a(0)=0$ and $a''<0$
already deliver $\alpha\geq 0$, equivalently $\varepsilon<1$, at every interior point. The one point at which $\alpha$
vanishes is the corner $x=0$, where $\alpha=a(0)-a'\!\left(0\right)\cdot 0=0$.}

\smalltitle{Allocation of capital.} The total capital stock is split between the two uses,
\begin{align}
K &= K^Y + K^R.
\label{eq:sp:setup:capital-allocation}
\end{align}

\smalltitle{Extraction and exploration.} Virgin materials are extracted from a stock $S$ of known reserves, which is augmented by exploration:
\begin{align}
C^N &= C^N(N,S,t),
&&
C^N_N>0,\quad C^N_{NN}\ge 0,\quad C^N_S<0, \quad C_t^N \leq 0, \quad C^N(0,S,t)=0
\label{eq:sp:setup:extraction-cost}
\\
C^D &= C^D(D,X, t),
&&
C^D_D>0,\quad C^D_{DD}\ge 0,\quad C^D_X>0, \quad C_t^D\leq 0, \quad C^D(0,X,t) =0.
\label{eq:sp:setup:discovery-cost}
\end{align}
In addition, we assume that discoveries are exhaustible in the sense that marginal exploration costs diverge at a finite cumulative level $\bar{X}_{\max}$: 
\begin{align}
    C^D_D(0,X,t)\rightarrow \infty \quad \text{as }X\rightarrow \bar{X}_{\max}<\infty. \label{eq:sp:setup:exhaustible-discovery}
\end{align}



\smalltitle{Waste collection and treatment.} All handled waste must be collected, and the treated share $\varpi$ of the handled flow carries an increasing marginal treatment cost:
\begin{align}
C^W &= \left[c^c(t)+c^T(\varpi,t)\right]H,
&&
c^T(0,t)=\left.\frac{dc^T}{d\varpi}\right|_{\varpi=0}=0,\quad
\frac{dc^T}{d\varpi}>0,\quad \frac{d^2c^T}{d\varpi^2}>0 \ \ \text{for }\varpi>0,
\label{eq:sp:setup:waste-cost}
\end{align}
and again both $c^c$ and $c^T$ are positive and decreasing (non-increasing) in time $t$. In addition, we impose the regularity assumption that marginal treatment costs are bounded:
\begin{align}
\left.\frac{dc^T}{d\varpi}\right|_{\varpi=1} &< \infty \qquad\text{for all }t .
\label{eq:sp:setup:cT-finite}
\end{align}
This assumption simplifies the model in the sense that the circularity limit is carried by the recycling technology alone.


\smalltitle{Accumulation.} The aggregate resource constraint and laws of motion for the remaining stocks are:
\begin{align}
Y &= C+I+\delta K+C^N+C^D+C^W, \label{eq:sp:setup:resource-constraint} \\
\dot K &= I, \qquad \dot S= D-N, \qquad \dot{X} = D \label{eq:sp:setup:reserves-motion} \\
\dot P &= \Xi-\theta(P)P, \qquad \theta'(P)<0,
\label{eq:sp:setup:pollution-motion}
\end{align}
together with the waste stock's law of motion~\eqref{eq:sp:setup:wastestock}, where the flow of waste that actually reaches the environment is\footnote{Making the stockpile environmentally free while it sits is a modelling choice, not a fact; a leakage term $\ell\mathcal W$ added to $\dot P$ would give landfills a flow cost alongside their asset value, and is recorded in Appendix~\ref{sec:app:extensions}.}
\begin{align}
\Xi &\equiv \left[1-\varpi+d^W\varpi\left(1-a\right)\right]H
\;=\; H\left(1-\varpi+d^W\varpi\right)-d^WR^R,
&&
0<d^W<1,
\label{eq:sp:setup:Xi}
\end{align}
the second equality using $a\varpi H=R^R$ from~\eqref{eq:sp:setup:recycling}.\footnote{Writing $\Xi$ this way is not cosmetic: it is linear in $H$ and $R^R$ separately, so every derivative of the pollution flow below is immediate, and no derivative of $a(\cdot)$ with respect to $\varpi$ or $H$ ever needs to be taken.}


\bigskip
In the following, where not explicitly useful, we ignore the time argument $t$ in various cost and production functions. 


\subsection{Materials accounting, process by process}
\label{sec:sp:setup:materialaccounting}
Every process below is described by a mass balance: what enters equals what leaves, distributed across waste and embodied material. The final good $Y$ is abstract in the sense that it is a composite of many different physical activities, but it is produced from physical raw material and is itself spent on physical or quasi-physical activities (consumption, capital formation, extraction effort, and so on). Tracking the materials balance principle correctly requires following the mass embodied in $Y$ through each of those uses.

Notationally, let $M^x$ and $W^x$ denote embodied materials and waste, respectively, from good/activity $x$. We let $\Phi^x$ and $\Omega^x$ denote embodied material and waste generation \emph{intensities}.

\smalltitle{Production.} 
Production uses raw material $R=N+R^R$ and capital $K^Y$, which depreciates at rate $\delta$. We assume that depreciation of capital returns embodied materials in $K^Y$ entirely as waste such that:
\begin{align}
    R+\Phi^K\delta K^Y &= W^Y+M^Y, &
W^Y &= \Omega^YY+\Phi^K\delta K^Y, &
M^Y &= \Phi^YY,
\label{eq:sp:setup:production}
\end{align}
where $\Phi^Y$ is the material intensity of the final good. The $\Phi^K\delta K^Y$ terms cancel on both sides, leaving
\begin{align}
R &= \Omega^YY+\Phi^YY.
\label{eq:sp:setup:production-reduced}
\end{align}


\smalltitle{Extraction and exploration.} Extraction uses the final good, through the cost function $C^N(N,S)$, to pull $N$ tonnes from the reserve $S$. Two physically distinct sources of waste are at play: the embodied material of the final good spent on extraction effort itself (machinery, energy, and so on, all ultimately final-good expenditure), and displaced material from the earth --- overburden and gangue brought up alongside the ore --- proportional to the tonnage extracted rather than to the effort expended, and never part of $Y$ to begin with:
\begin{align}
W^N &= \Omega^{N,Y}C^N(N,S)+\Omega^{N,S}N.
\label{eq:sp:setup:extraction}
\end{align}
The same logic applies to exploration, which spends $C^D(D,X)$ of the final good to add $D$ to the stock of discoveries:
\begin{align}
W^D &= \Omega^{D,Y}C^D(D,X)+\Omega^{D,S}D.
\label{eq:sp:setup:exploration}
\end{align}

Here, the $\Omega^{N,Y}C^N+\Omega^{D,Y}C^D$ represent embodied materials from $Y$ that are diffused as waste materials after being used in exploration and extraction processes, while $\Omega^{N,S} N $ and $\Omega^{D,S} D$ are nature-attributable components that represent additional mass that never passes through $Y$ and thus never draws on $R$. 

\textit{Note:} The waste generation that comes directly from exploration and extraction should ideally also be extracted from some nature stock so the materials are not coming "like manna from heaven". For now, however, we ignore this stock-flow accounting assuming that this will not be the limiting factor for the economy even in the long run.




\smalltitle{Waste collection and treatment, and consumption.} Both activities draw the final good and return its entire embodied content as waste, since neither stores materials durably: consumption goods are used up, and treating waste is itself an activity that consumes final good and leaves nothing durable behind. Consumption is already measured in final-good units, so $\Omega^C$ needs no further translation. Waste collection and treatment is not: the physical quantity it processes is the handled flow $H$, a tonnage, while the final-good resources it uses are $C^W=\left[c^c+c^T(\varpi)\right]H$ from~\eqref{eq:sp:setup:waste-cost}. 


Waste generation is proportional to the economic activity, representing the diffusion of embodied materials in $Y$ as waste materials. Thus, we simply define waste generation streams:\footnote{We could in principle also assume that there are additional waste generation from collection and treatment of $W$, but the only story I can think of here is that there may be residual waste generated from the recycling process and this story is easier to simply subsume in the recycling function instead of creating a new waste flow.}
\begin{align}
    W^W = \Omega^{W} C^W \qquad W^C= \Omega^C C, \label{eq:sp:setup:wc}
\end{align}


\smalltitle{Investment.} Gross resources devoted to capital formation are
\begin{align}
G &\equiv I+\delta K,
\label{eq:sp:setup:gross-formation}
\end{align}
net accumulation plus replacement of depreciated capital, matching the $I$ and $\delta K$ terms in~\eqref{eq:sp:setup:resource-constraint}. We assume that embodied materials from final goods used in accumulating durables are entirely embodied in the capital, and the waste generation is then associated with the depreciation. This implies
\begin{align}
    W^I &= 0, & M^I &= \Phi^IG.
\label{eq:sp:setup:investment}
\end{align}


\smalltitle{Recycling.} Recycling uses capital $K^R$ and waste materials only, drawing no final good directly; its own depreciation releases waste in the same way as $K^Y$'s:
\begin{align}
W^R &= \Phi^K\delta K^R, & R^R &= a\!\left(x\right)\varpi H.
\label{eq:sp:setup:recycling-mass}
\end{align}
We note that recycling in itself may produce residual waste. However, in our model with a single waste material type, we interpret the recycling function as including the subsequent treatment of the residual waste, i.e. measuring a net recycling efficiency taking residual waste generation into account. 

\subsection{The embodied-material stock}
\label{sec:sp:setup:embodied-stock}
Because gross investment is embodied at the current intensity $\Phi^I(t)$ while depreciation releases material at the stock's average intensity $\Phi^K(t)$, the total mass of material currently embodied in the capital stock, $M^K\equiv\Phi^KK$, is itself a state variable:
\begin{align}
\dot M^K &= \Phi^I(t)\left(I+\delta K\right)-\delta M^K,
&
M^K(0) \text{ given},
\label{eq:sp:setup:Phi-motion}
\end{align}
structurally identical to $K$'s own law of motion $\dot K=I$, with gross investment $I+\delta K$ replacing $I$ and the flow intensity $\Phi^I(t)$ replacing the constant $1$. Equivalently, in ratio form,
\begin{align}
\dot\Phi^K &= \left[\Phi^I(t)-\Phi^K\right]\frac{I+\delta K}{K},
\label{eq:sp:setup:phiK-motion}
\end{align}
so $\Phi^K(t)$ is an investment-weighted moving average of past values of $\Phi^I$: it drifts toward the current vintage's intensity at a rate set by the gross investment rate, and is stationary exactly when $\Phi^I(t)=\Phi^K(t)$.


\subsection{Aggregate material balances}
\label{sec:sp:setup:aggregatematerial}
Tracking the embodied materials in the final goods, the material balance principle requires that
\begin{align}
    \Phi^Y Y = R - \Omega^Y Y = \Omega^{N,Y}C^N + \Omega^{D,Y}C^D + \Omega^C C + \Omega^W C^W +\Phi^I G, \label{eq:sp:setup:ybalance_generic}
\end{align}
where the first equality comes from \eqref{eq:sp:setup:production-reduced} and the second comes from summing over all the processes that the final goods are eventually used (extraction, exploration, consumption, investments, and treatment). We can also rearrange the last equality and instead write
\begin{align}
    R =  \Omega^Y Y + \Omega^{N,Y}C^N + \Omega^{D,Y}C^D + \Omega^C C + \Omega^W C^W + \Phi^I G \label{eq:sp:setup:ybalance_generic2},
\end{align}
showing that the new materials that goes into production either ends up as waste or embodied in new investment goods. Furthermore, total waste generation $W$ is given by summing over all the relevant sources:
\begin{align}
    W =& \Omega^Y Y + \Omega^{N,Y}C^N + \Omega^{D,Y}C^D + \Omega^C C + \Omega^W C^W + \Omega^{N,S}N + \Omega^{D,S}D + \Phi^K \delta K \notag \\
    =& R - \Phi^I G + \Omega^{N,S}N + \Omega^{D,S}D + \Phi^K\delta K,
    \label{eq:sp:setup:totalwaste_generic}
\end{align}
where the second line uses \eqref{eq:sp:setup:ybalance_generic2}. 



\bigskip
How we treat parameters $\Omega, \Phi$ from here on out, however, is essential for the interpretation of the model and how the material balance principles constraint the model. The first condition \eqref{eq:sp:setup:ybalance_generic} essentially identifies $\Phi^Y\geq0$, but depending on if other intensities are fixed or endogenous, this can either be a simple definition of $\Phi^Y$ or a real technological constraint: If all other $\Phi, \Omega$ parameters are exogenous, for instance, this severely limits what $Y$ are actually feasible (as $\Phi^Y \geq 0$ has to hold). If, alternatively, the levels of coefficients on the right-hand side scale with the material inputs, this is just a way to compute $\Phi^Y$ without restricting $Y$. 




We generally do not think that there are any good reason for having exogenous levels of waste generating intensities from using final goods. Thus, we refactor using $\Omega^j = \Omega \omega^j$ for all the streams tied to embodied materials in $Y$, reflecting that the waste generation (in tonnes) depends on the tonnes of materials used in final goods production. In this case, we can rewrite \eqref{eq:sp:setup:ybalance_generic2} as
\begin{align}
    R =  \Omega \left(\omega^Y Y + \omega^{N}C^N + \omega^{D}C^D + \omega^C C + \omega^W C^W\right) + \Phi^I G \label{eq:sp:setup:ybalance_Omega},
\end{align}
using this to define $\Omega$ (that adjusts endogenously), and the only real constraint is then that $\Omega\geq0$.

It is worth being explicit about what this layer of the accounting is for. The allocation is driven by a small set of physical objects (extraction, recycled material, the material intensity of new capital, depreciation, the treatment share and the recycling function) whereas material-flow statistics are reported as \emph{intensities}: waste or material input per unit of economic activity. $\Omega$ and the relative intensities $\omega^j$ are the bridge between the two. They are how the model's mass flows are mapped into, and calibrated against, published accounts, and that remains their role even at points where they turn out not to affect the allocation at all.

Finally, the material intensity of new capital goods, $\Phi^I$, can be interpreted in different ways: We can assume that it requires a very specific level of materials to produce durables, which then requires $\Phi^I$ to be exogenous in levels, we can even assume that this level is constant over time (simplifying the model by removing a state variable), or we may assume that the material intensity of durables follows that of the final good (in levels at least). The documentation uses this last assumption and assumes $\Phi^I = \Omega \phi^I$ and include the other two versions in appendices \ref{sec:app:exogenousPhiI} and \ref{sec:app:constantPhiI}. Under this assumption we have:
\begin{align}
    R =  \Omega \left(\omega^Y Y + \omega^{N}C^N + \omega^{D}C^D + \omega^C C + \omega^W C^W+\phi^I G\right) \label{eq:sp:setup:ybalance_OmegaPhi},
\end{align}
where the condition $\Omega\geq0$ now holds per construction.

\smalltitle{Aggregate balances and circularity multiplier.}
Next, to simplify some of the equations a bit (at the expense of more notation though), we define the waste bases
\begin{align}
    \mathcal D &\equiv \omega^YY+\omega^{N}C^N(N,S)+\omega^{D}C^D(D,X)+\omega^WC^W+\omega^CC,
&
\mathcal N &\equiv \Omega^{N,S}N+\Omega^{D,S}D,
\label{eq:sp:setup:base}
\end{align}
with $\mathcal D$ collecting flows measured in final-good units, carrying embodied material drawn through $R$, while $\mathcal N$ gathers the two physical flows that do not, and with $C^W=c^WH$, $c^W\equiv c^c+c^T(\varpi)$ throughout. Given this additional notation, our first "result" comes from substituting~\eqref{eq:sp:setup:ybalance_generic2} into the waste sum~\eqref{eq:sp:setup:totalwaste_generic} to get:\footnote{The substitution and sums cancels $\Omega$ and every $\omega^j$ from the accounted part, and~\eqref{eq:sp:setup:Phi-motion} collects the two remaining capital terms into $\dot M^K$.}
\begin{align}
\boxed{\;W \;=\; R-\Phi^IG+\delta M^K+\mathcal N \;=\; R-\dot M^K+\mathcal N\;}
\label{eq:sp:setup:ledger}
\end{align}
\emph{Waste generated equals the material entering production, less the net accumulation of material embodied in the capital stock, plus the mass displaced from nature that never entered production at all.} That is the entire physical content of the materials balance principle for this economy, and it is what the planner's problem below is constrained by. One feature is worth marking: the average intensity $\Phi^K$ appears nowhere in it, so it is needed only to report the composition of the capital stock, never to state the balance.


Next, substituting $R=N+a\varpi H$ into~\eqref{eq:sp:setup:ledger} and using $W=\dot{\mathcal W}+H$ from~\eqref{eq:sp:setup:wastestock} therefore gives
\begin{align}
\left(1-a\varpi\right)H &= N-\dot M^K-\dot{\mathcal W}+\mathcal N,
&
a\varpi&\le\bar a\le 1 .
\label{eq:sp:setup:ledger-collected}
\end{align}
The left-hand side is the \emph{leakage flow}: of every handled tonne, the fraction $1-a\varpi$ leaves circulation for good, and it must be covered by the sources on the right --- virgin extraction and displaced mass, net of what the economy accumulates in its two anthropogenic stocks, $M^K$ and $\mathcal W$.

The factor $\left(1-a\varpi\right)^{-1}$ is the \emph{circularity multiplier}: the number of times a tonne of material passes through production before it finally leaves the economy. With the waste stock in place it is a steady-state object rather than an instantaneous identity: whenever the material block is stationary ($\dot{\mathcal W}=\dot M^K=0$, so that $H=W$) and no mass is displaced ($\mathcal N=0$), \eqref{eq:sp:setup:ledger-collected} reads $N=W\left(1-a\varpi\right)$ --- virgin extraction equals waste not recycled --- and each tonne entering the economy serves production $\left(1-a\varpi\right)^{-1}$ times on its way through. A fully circular economy --- one that extracts nothing while processing a positive throughput --- requires $a\varpi=1$, and whether that is out of reach is exactly the content of the yield-ceiling case distinction.

Under the hard ceiling~(H), $a\varpi\le\bar a<1$, so the multiplier is bounded by $\left(1-\bar a\right)^{-1}$ uniformly: a tonne entering the economy serves production at most $\left(1-\bar a\right)^{-1}$ times before it is lost, whatever the path. Circularity can stretch the economy's material budget, never unbound it.

Under the soft ceiling~(S) no such uniform bound exists. The multiplier is finite at every date --- $a\left(x,t\right)<1$ for every finite $x$, and Assumption~\eqref{eq:sp:setup:cT-finite} does not help, since $\varpi\le1$ --- but it is unbounded along paths on which recycling capital per treated tonne diverges, so full circularity is approached rather than attained. What disciplines throughput in this case is the waste stock itself. Recovered material is drawn from the handled flow, so at every date and under either ceiling
\begin{align}
R^R &= a\varpi\,\mu\mathcal W \;\le\; \bar a\,\mu\mathcal W,
&
R &\le N+\bar a\,\mu\mathcal W .
\label{eq:sp:setup:throughput-cap}
\end{align}
\emph{Circular throughput is bounded by the retained stock of waste times the handling rate.} However close $a\varpi$ comes to one, a closed loop cannot manufacture throughput out of nothing: sustaining a recycled supply $R^R$ requires holding an anthropogenic reserve of at least $R^R/\left(\bar a\mu\right)$ tonnes. The multiplier measures how far a tonne stretches over its lifetime in the economy; the stock and the handling rate cap how fast material can circulate at any moment. The economically meaningful long-run question under the soft ceiling is then whether the leakage the multiplier measures is cumulatively finite, and Section~\ref{sec:sp:longrun} works with the form~\eqref{eq:sp:setup:ledger-collected} throughout, since it is the one statement of the ledger that does not presume which ceiling case holds.



% With $\Phi^I = \Omega \phi^I$, material and waste intensities of final goods $(\Omega, \Phi^Y)$ are both non-negative per construction. Rearranging the balance condition in \eqref{eq:sp:setup:ybalance_OmegaPhi} and \eqref{eq:sp:setup:ybalance_generic} gives
% \begin{align}
% \Omega &= \frac{R}{\mathcal D+\phi^IG}\geq0, \qquad \Phi^Y = \frac{\Omega\left(\mathcal D + \phi^I G\right) - \Omega^Y Y}{Y}\geq0, 
% \label{eq:sp:setup:Omega-phi}
% \end{align}
% where $\Omega\geq0$ comes from all parts being non-negative and $\Phi^Y\geq0$ comes from $\Omega \mathcal{D}\geq \Omega^Y Y$.



